\subsection{Glossario}
Per evitare fraintendimenti legati al linguaggio impiegato nella documentazione, viene predisposto un Glossario in cui ogni voce è corredata di una spiegazione volta a chiarirne il significato. Tutti i termini tecnici, gli acronimi e le parole suscettibili di equivoci riportati nei documenti principali saranno:

\begin{itemize}
  \item resi in \textit{corsivo};
  \item accompagnati dalla lettera “G” in pedice (ad es.\ \textit{Termine}$_G$);
  \item collegati, tramite link, alla rispettiva voce del Glossario.
\end{itemize}

In questo modo, ogni volta che un termine definito compare nel testo, sarà immediatamente riconoscibile e rimanderà alla definizione ufficiale, eliminando così qualsiasi ambiguità.